\documentclass[12pt]{article}
\usepackage[utf8]{inputenc}
\usepackage[T1]{fontenc}
\usepackage{hyperref}
\title{Hadrami Language Resources for K-12 Curriculum Development}
\author{Compiled by Claude Sonnet 5.5 (claude-sonnet-5-5) for Andr\'as Kornai}
\date{}
\begin{document}
\maketitle

\section{Introduction}
Hadrami, also called Hadramautic (ISO 639-3 xhd, Glottocode hadr1235), is an extinct Old South Arabian Semitic language. It was used from roughly 800 BC to 600 AD in the Kingdom of Hadhramaut, around the capital Shabwa in present-day Yemen. It was also used at the trading post Sumhuram (Khor Rori, Oman). Nearly all evidence comes from inscriptions in the Ancient South Arabian script. After the Himyarite conquest, Sabaic became the official language. The people of the region now mainly speak Hadhrami Arabic, which is a separate language. No community of heritage speakers exists, and no school materials were found.

\section{Primary Educational Resources}
I found no primers, textbooks or school programmes for ancient Hadramautic.

\section{Linguistic and Reference Materials}
\begin{itemize}
\item Wikipedia, Hadramautic language: \url{https://en.wikipedia.org/wiki/Hadramautic_language}
\item Wikipedia, Kingdom of Hadhramaut (historical context): \url{https://en.wikipedia.org/wiki/Kingdom_of_Hadhramaut}
\item Wiktionary, category of Hadrami language entries: \url{https://en.wiktionary.org/wiki/Category:Hadrami_language}
\item DASI, Digital Archive for the Study of pre-Islamic Arabian Inscriptions, an open TEI-Epidoc inscription archive covering Hadramitic along with Sabaic, Minaic and Qatabanic: \url{https://dasi.cnr.it/index.php?id=83}
\end{itemize}

\section{Religious and Community Resources}
I found no religious or community resources specific to Hadramautic. Hadhrami Arabic, the modern spoken language of the region, is a different language (see Wikipedia: \url{https://wikipedia.com/wiki/ISO_639:ayh}).

\section{Digital and Audio Resources}
DASI (above) is the main digital resource. I found no audio or video resources.

\section{Curriculum Development Context}
As a language known only from ancient inscriptions, Hadramautic suits a history, script and cultural-heritage unit more than a language-revitalization curriculum. Possible content includes the Ancient South Arabian script, the incense-trade kingdom, and a few inscription texts read with translations from DASI. The sources give slightly different dates and name variants, so a curriculum should check them against scholarly editions.

\end{document}